\chapter{Lebesgue 测度}
\label{chap:lebesgue-measure}

长度应满足三项要求：区间的长度是端点之差，平移不改变长度，可数个互不
相交部分的长度可以相加。Jordan 理论只允许有限覆盖，因而稠密可数集与
整个区间无法区分。把覆盖放宽为可数覆盖，再应用上一章的延拓过程，得到
Lebesgue 测度。

\section{实直线上的外长度}

对任意 $E\subseteq\R$，定义
\begin{equation}
 m^*(E)=\inf\left\{\sum_{n=1}^{\infty}|I_n|:
 E\subseteq\bigcup_{n=1}^{\infty}I_n, I_n\text{ 为开区间}\right\}.
 \label{eq:lebesgue-outer-measure-line}
\end{equation}
覆盖族中允许空区间，其长度为零；由覆盖外测度的一般构造，$m^*$ 是外测度。
允许无界区间不会改变定义；
需要有限总代价时，所用区间必可取为有界区间。

\begin{lemma}[区间的外长度]
\label{lem:interval-outer-length}
若 $a<b$，则
\[
 m^*([a,b])=m^*((a,b))=b-a.
\]
同样结论适用于半开区间。
\end{lemma}

\begin{proof}
区间 $(a-\varepsilon,b+\varepsilon)$ 给出上界，令
$\varepsilon\downarrow0$。为证下界，设开区间列覆盖 $[a,b]$。紧致性
给出有限子覆盖。有限开区间覆盖闭区间时，其长度和至少为 $b-a$：删去
无用区间并按左端点排列，从含 $a$ 的区间开始，每次选取在已覆盖右端点
之前起始且右端点最远的区间，所得有限链的长度和至少为首末端点之差。
因此任意可数覆盖的总代价均不小于 $b-a$。

开区间包含任意闭子区间 $[a+\varepsilon,b-\varepsilon]$，故其外长度至少
$b-a-2\varepsilon$；令 $\varepsilon\downarrow0$。端点是外测度零集，
所以半开区间相同。
\end{proof}

平移覆盖中的每个区间可知
\begin{equation}
 m^*(E+t)=m^*(E),
 \label{eq:outer-measure-translation}
\end{equation}
反向平移给出等号。若集合 $A,B$ 的距离为正，记
$d(A,B)=\delta>0$。外长度定义中的区间可以限定为长度小于 $\delta$ 的
区间：把每个有界覆盖区间作有限细分，并将端点处略微重叠，增加的总长度
可以预先控制为任意小。这样的短区间不可能同时遇到 $A$ 与 $B$。因此一组
短区间覆盖 $A\cup B$ 时，分别取其中遇到 $A$ 和遇到 $B$ 的区间，便得到
$A$、$B$ 的两个覆盖。对覆盖代价取下确界可得
\begin{equation}
 m^*(A\cup B)=m^*(A)+m^*(B).
 \label{eq:metric-outer-additivity}
\end{equation}
这种外测度称为度量外测度。

\section{Lebesgue 可测集}

\begin{definition}
$m^*$ 的 Carathéodory 可测集称为\term{Lebesgue 可测集}，其全体记为
$\mathcal L(\R)$；限制测度记为 $m$，称为\term{Lebesgue 测度}。
\end{definition}

度量外测度的正距离可加性说明所有闭集可测，因而所有 Borel 集可测。由
Carathéodory 定理，$m$ 完备；所以
\[
 \mathcal B(\R)\subseteq\mathcal L(\R),
\]
并且每个 Borel 零测集的任意子集都 Lebesgue 可测。后一类子集未必是
Borel 集，这正是两个 $\sigma$-代数的差别来源。

区间长度、平移不变性和可数可加性现在同时成立。特别地，每个单点零测，
故每个至多可数集零测。与 Jordan 零测不同，覆盖中允许可数多个区间，
所以 $\Q\cap[0,1]$ 的 Lebesgue 测度为零。

\begin{theorem}[一维 Lebesgue 测度的正则性]
\label{thm:lebesgue-regularity-line}
对每个 $E\in\mathcal L(\R)$ 和 $\varepsilon>0$，存在开集 $O\supseteq E$
使
\[
 m(O\setminus E)<\varepsilon.
\]
因此
\[
 m(E)=\inf\{m(O):E\subseteq O, O\text{ 开}\}.
\]
若 $m(E)<\infty$，则
\[
 m(E)=\sup\{m(K):K\subseteq E, K\text{ 紧}\}.
\]
\end{theorem}

\begin{proof}
先设 $m(E)<\infty$。由外测度定义取开区间覆盖，使长度和小于
$m(E)+\varepsilon$；其并 $O$ 开，且有限性允许相减，从而
$m(O\setminus E)<\varepsilon$。

一般地，把整数半开区间依次编号为 $I_j$，并令 $E_j=E\cap I_j$。每个
$E_j$ 可测且测度有限，并且 $E=\bigcup_jE_j$。选正数 $c_j$ 使
$\sum_jc_j=1$，再为每个 $E_j$ 取开集 $O_j\supseteq E_j$，满足
$m(O_j\setminus E_j)<\varepsilon c_j$。于是 $O=\bigcup_jO_j$ 开，
$E\subseteq O$，而
\[
 O\setminus E\subseteq\bigcup_j(O_j\setminus E_j),
 \qquad m(O\setminus E)<\varepsilon.
\]
这也证明了外正则性公式，包括 $m(E)=\infty$ 的情形。

以下仍设 $m(E)<\infty$。由
$E\cap[-N,N]\uparrow E$，可选 $N$ 使
$m(E\setminus[-N,N])<\varepsilon/2$。对有限测度集
$F=[-N,N]\setminus E$ 取开外包 $O$，使
$m(O\setminus([-N,N]\setminus E))<\varepsilon/2$。令
\[
 K=[-N,N]\setminus O.
\]
则 $K$ 紧且 $K\subseteq E$。并且
\[
 E\setminus K
 \subseteq (E\setminus[-N,N])\cup(O\setminus F),
\]
所以 $m(E\setminus K)<\varepsilon$。这给出内正则性公式。
\end{proof}

正则性还给出常用的 Borel 代表。对每个 $n$ 取开集
$O_n\supseteq E$，使 $m(O_n\setminus E)<2^{-n}$，则
$G=\bigcap_nO_n$ 是 $G_\delta$ 集，$E\subseteq G$ 且
$m(G\setminus E)=0$。若 $m(E)<\infty$，取紧集
$K_n\subseteq E$，使 $m(E\setminus K_n)<2^{-n}$，则
$F=\bigcup_nK_n$ 是 $F_\sigma$ 集且 $m(E\setminus F)=0$。一般可测集先
按整数区间分成有限测度块，再对各块作这一构造；可数个 $F_\sigma$ 集的并
仍是 $F_\sigma$ 集。因此每个 Lebesgue 可测集都与某个 $G_\delta$ 集、
某个 $F_\sigma$ 集分别相差零测集。

\section{欧氏空间中的 Lebesgue 测度}

设 $d\ge1$ 为整数。对有界开、闭或半开长方体，规定体积为各边长之积，
空集体积为零。有限个半开长方体的不交并，其体积按各块相加；两种有限
分解可用所有坐标端点作共同网格细分，因此结果与分解无关。不过这些集合
不包含全空间，不能直接套用上一章在集合代数上陈述的延拓定理。
以下从开长方体覆盖构造。
\[
 m_d^*(E)=\inf\left\{\sum_{n\ge1}v(O_n):
 E\subseteq\bigcup_{n\ge1}O_n,\ O_n\text{ 为有界开长方体或空集}\right\}.
\]
它由覆盖构造成为外测度。

\begin{proposition}
每个有界开、闭或半开长方体 $R$ 满足 $m_d^*(R)=v(R)$。
\end{proposition}
\begin{proof}
稍向外扩张各面，得到覆盖 $R$ 的开长方体，其体积任意接近 $v(R)$，
故 $m_d^*(R)\le v(R)$。若 $R$ 非退化，给定 $\varepsilon>0$，稍向内收缩
各面，可取紧长方体 $K\subset R$ 且 $v(K)>v(R)-\varepsilon$。
任意开长方体覆盖 $(O_n)$ 覆盖 $K$，由紧致性存在有限子覆盖。
按这有限族及 $K$ 的全部坐标端点细分，逐格比较体积，得到
\[
 v(K)\le\sum_{j=1}^s v(O_{n_j})\le\sum_{n\ge1}v(O_n).
\]
网格面不影响有限体积和；这也是前文有限长方体覆盖的体积不等式。
令 $\varepsilon\downarrow0$，再对覆盖取下确界，得反向不等式。
退化长方体可以装入体积任意小的开长方体，故其外测度为零。
\end{proof}

任给 $\delta>0$，覆盖可限定为直径小于 $\delta$ 的开长方体。
每个有界覆盖长方体作有限细分，再微量扩大各块，就能开覆盖原长方体，
保持直径小于 $\delta$，并使第 $n$ 个长方体增加的总体积小于
$\varepsilon2^{-n}$。因此外测度有限时，总增加量可任意小。
距离为正的两个集合不能同时与这样的短长方体相交，分开覆盖后取下确界，
得到度量可加性；若并集外测度无穷，所需反向不等式自动成立。
由上一章度量外测度的证明，所有闭集、从而所有Borel集可测。

定义 $\mathcal L(\R^d)=\mathcal M_{m_d^*}$，以及
$m_d=m_d^*|_{\mathcal L(\R^d)}$。Carathéodory 定理给出完备测度，
它在长方体上等于体积；整数单位半开立方体覆盖全空间，故它是
$\sigma$-有限的。不致混淆时仍写 $m$ 或 $|E|$。

它在Borel集上的限制由长方体体积唯一确定。固定 $Q_N=[-N,N)^d$，
其中有限个半开长方体的并构成 $Q_N$ 上的相对代数，总测度有限。
欧氏开集有可数有理长方体基，而开长方体又可写成内部半开长方体的可数并，
故这个相对代数生成 $\mathcal B(Q_N)$。两个给长方体相同体积的Borel测度，
由上一章有限延拓唯一性在 $Q_N$ 上相等；令 $N\uparrow\infty$ 即在全部
Borel集上相等。这里所用代数始终包含其相对全集。

平移及非零伸缩把开长方体覆盖一一对应，故对任意集合 $E$，
\[
 m_d^*(x+E)=m_d^*(E),\qquad
 m_d^*(rE)=|r|^d m_d^*(E)\quad(r\ne0).
\]
将任意测试集作逆变换，在Carathéodory等式两边同乘相同的有限正数，
可知这些变换保持Lebesgue可测性。因此对可测集，
\[
 m_d(x+E)=m_d(E),\qquad m_d(rE)=|r|^d m_d(E)\quad(r\ne0).
\]
若 $r=0$，像为空集或单点，测度为零；当 $m_d(E)=\infty$ 时，
不把它写成未定义的乘积 $0\cdot\infty$。

正则性也由覆盖定义得到。设 $m_d(E)<\infty$，选取总体积小于
$m_d(E)+\varepsilon$ 的开长方体覆盖。其并为开外包 $O$，并且
$m_d(O\setminus E)<\varepsilon$。一般 $E$ 先与整数单位半开立方体逐个相交，
对第 $j$ 块使用误差 $\varepsilon2^{-j}$（$j\ge1$），再取开外包的并。
由于多出的部分包含于各块多出部分之并，仍有差集测度小于 $\varepsilon$。

若 $m_d(E)<\infty$，取大闭立方体 $Q$ 使
$m_d(E\setminus Q)<\varepsilon/2$。对 $Q\setminus E$ 取开外包 $O$，
额外测度小于 $\varepsilon/2$，则 $K=Q\setminus O\subseteq E$ 紧，且
$m_d(E\setminus K)<\varepsilon$。这证明了开外逼近和紧内逼近，
不需要预先分解所有开集。

由开外包逐次取交，每个可测 $E$ 都包含于某个Borel集 $G$，且
$m_d(G\setminus E)=0$。再对这个零测差集取开外包的交，得到包含它的
Borel零测集 $H$。于是 $E\mathbin\triangle G\subseteq H$，说明
$\mathcal L(\R^d)$ 恰是Borel测度的完备化；反向包含来自Carathéodory
测度的完备性。覆盖构造与Borel测度完备化因而给出同一测度空间。

\section{与 Jordan 测度的关系}

\begin{theorem}[Jordan 与 Lebesgue 测度]
\label{thm:jordan-lebesgue-comparison}
若 $E\subset\R^d$ 有界且 Jordan 可测，则 $E$ Lebesgue 可测，并且其
Jordan 体积等于 $m_d(E)$。反过来，有界集 $E$ Jordan 可测当且仅当
$m_d(\partial E)=0$。
\end{theorem}

\begin{proof}
定理 \ref{thm:jordan-indicator-boundary} 已经证明：有界集 $E$ 的指标函数
Riemann可积，当且仅当 $\partial E$ 为Jordan零测集。这正是Jordan可测性
的有限网格判据。该证明把网格所有面另作零测覆盖；只覆盖指标函数振幅非零
的网格块不足以覆盖全部边界，尤其不能遗漏外围矩形的面。

若 $E$ Jordan 可测，上述有限覆盖也是 Lebesgue 外测度的允许覆盖，故
$m_d(\partial E)=0$。集合 $E^\circ$ 可测，且
\[
 E^\circ\subseteq E\subseteq\overline E
 =E^\circ\cup\partial E.
\]
Lebesgue 测度的完备性于是给出 $E$ 可测，并且
$m_d(E)=m_d(E^\circ)=m_d(\overline E)$。对外和、内和对应的有限长方体并
应用测度的单调性，得到
\[
 \text{内和}\le m_d(E)\le\text{外和}.
\]
当分割使两者之差趋于零时，它们的共同极限既是 Jordan 体积，也是
$m_d(E)$。

反过来，若 $m_d(\partial E)=0$，则紧集 $\partial E$ 有可数开长方体覆盖，
其总体积可任意小。紧致性从中抽出有限子覆盖，且所取有限族的总体积仍小于
原可数族的总体积。因此 $\partial E$ 是 Jordan 零测集，由有限网格判据知
$E$ Jordan 可测。
\end{proof}

Lebesgue 可测有界集未必 Jordan 可测。例如 $\Q\cap[0,1]$ 可测且零测，
但边界是整个 $[0,1]$。Lebesgue 理论允许可数覆盖及零测修改，因而比
Jordan 理论稳定得多。

\section*{习题}
\addcontentsline{toc}{section}{习题}

\subsection*{A组　外长度与可测性}
\begin{enumerate}[label=\textbf{\thechapter.\arabic*.}]
 \exerciseitem{interval-cover-length-lemma}
 详细证明有限开区间覆盖 $[a,b]$ 时长度和至少为 $b-a$。
 \exerciseitem{lebesgue-singletons-countable}
 证明单点和至多可数集的 Lebesgue 测度为零。
 \exerciseitem{translation-invariance-outer}
 由覆盖定义证明 $m^*(E+t)=m^*(E)$。
 \exerciseitem{scaling-lebesgue-measure}
 证明 $m_d(rE)=|r|^dm_d(E)$，并处理 $r=0$。
 \exerciseitem{open-set-countable-intervals}
 证明 $\R$ 中每个开集是至多可数个两两不交开区间的并。
\end{enumerate}

\subsection*{B组　正则性与 Borel 集}
\begin{enumerate}[label=\textbf{\thechapter.\arabic*.},resume]
 \exerciseitem{lebesgue-gdelta-envelope}
 证明每个 Lebesgue 可测集 $E$ 包含于某个 $G_\delta$ 集 $G$，且
 $m(G\setminus E)=0$。
 \exerciseitem{lebesgue-fsigma-kernel}
 证明每个有限测度 Lebesgue 可测集含有 $F_\sigma$ 集 $F$，且
 $m(E\setminus F)=0$。
 \exerciseitem{borel-versus-lebesgue-cardinality}
 设 $C=\{\sum_{n\ge1}2\varepsilon_n3^{-n}:\varepsilon_n\in\{0,1\}\}$。
 证明 $C$ 是闭零测集且与 $\mathcal P(\Npos)$ 等势，并由此证明
 $\mathcal B(\R)\ne\mathcal L(\R)$。为估计Borel集族的大小，可把可数并和
 补运算编码为 $\N$ 的有限序列组成的可数树；须说明这种编码为何涵盖全部Borel集。
 \exerciseitem{regular-approximation-symmetric-difference}
 若 $E$ 可测且有限测度，证明对每个 $\varepsilon>0$ 存在紧集 $K$ 与开集
 $O$，使 $K\subseteq E\subseteq O$ 且 $m(O\setminus K)<\varepsilon$。
 \exerciseitem{boundary-null-graph}
 证明紧矩形上 Lipschitz 函数的图像在高一维 Euclidean 空间中 Lebesgue
 零测。
\end{enumerate}

\subsection*{C组　Jordan 与 Lebesgue 理论}
\begin{enumerate}[label=\textbf{\thechapter.\arabic*.},resume]
 \exerciseitem{jordan-lebesgue-volume-equality}
 用网格内外逼近完整证明定理 \ref{thm:jordan-lebesgue-comparison} 中体积
 相等的部分。
 \exerciseitem{rationals-not-jordan-measurable}
 证明 $\Q\cap[0,1]$ Lebesgue 可测但不 Jordan 可测。
 \exerciseitem{boundary-positive-fat-set-preview}
 构造或引用一个正测度、无内点的闭集 $F\subset[0,1]$，说明 $F$ 不
 Jordan 可测。
 \exerciseitem{measure-zero-image-linear}
 若 $A:\R^d\to\R^k$ 为线性映射且 $N\subset\R^d$ 零测，证明在
 $k=d$ 时 $A(N)$ 零测；分析 $k<d$ 时应如何理解结论。
 \optionalexerciseitem{uniqueness-translation-invariant-measure}
 设 $\nu$ 是 $\R$ 上的平移不变 Borel 测度，有限于有界区间且
 $\nu([0,1))=1$。证明 $\nu$ 等于 Lebesgue 测度。
\end{enumerate}
